\section{Visão Geral}

\hspace{0.5cm} A arquitetura do sistema ``Raízes do Nordeste'' segue o padrão de camadas,
separando responsabilidades em níveis distintos: apresentação via rotas, lógica de negócio em serviços
e persistência em repositórios.

Cada domínio segue essa estrutura de forma independente entre acesso a dados no repositório,
regras de negócio no serviço e contrato HTTP nas rotas, facilitando evolução isolada de cada feature, testes unitários e refatorações sem impactos em cascata.
Quanto à infraestrutura, é utilizado um banco de dados relacional com tabelas normalizadas com eficiência nas consultas.

Por fim, a estrutura de pastas do projeto reflete essa organização em camadas, com diretórios dedicados para cada domínio,
além de uma seção de infraestrutura para utilitários, configuração de banco e migrações.

\begin{verbatim}
Domínios:
├── Audit
│   ├── audit_decorator.py
│   ├── audit_middleware.py
│   ├── audit_repository.py
│   ├── audit_service.py
│   └── audit.py
├── Auth/
│   ├── auth.py
│   ├── auth_service.py
│   └── auth_routers.py
├── Catalogue
│   ├── catalogue_repository.py
│   ├── catalogue_routers.py
│   ├── catalogue_service.py
│   ├── catalogue.py
│   └── dish_ingredient.py
├── Costumer/
├── Employee/
├── Inventory/
└── Order/
Infraestrutura:
├── Utils/
│   ├── base_repository.py
│   ├── base_service.py
│   ├── update_helper.py
│   ├── ownership_decorator.py
│   ├── address.py
│   └── validations.py
├── Database/
│   └── db_config.py
├── migrations/
├── main.py
└── alembic.ini
\end{verbatim}
